% =============================================================================
% CONCLUSION - CITA Paper in ALKALI Format
% =============================================================================

\vspace{-2mm}
\section{Conclusion}
\label{sec:conclusion}

This paper introduces \textbf{CITA (Contrastive Instruction-Tuned Alignment)}, a framework enabling instruction-aware alignment in large language models. Unlike DPO~\cite{rafailov2023direct} and RLHF~\cite{ouyang2022training} which embed static alignment into weights, CITA allows models to dynamically adapt behavior based on natural language instructions at inference time.

\vspace{-2mm}
\subsection{Key Findings}
\label{sec:key_findings}

\begin{enumerate}[leftmargin=*,itemsep=2pt]
    \item \textbf{CITA outperforms DPO in instruction-awareness}: Measuring NoInstruct $\rightarrow$ Instruct improvement, CITA wins \textbf{3/5 benchmarks} (TruthfulQA, Length Control, AQI), while DPO wins 2/5 (ISD, Conditional Safety).

    \item \textbf{Correct directional adaptation}: On TruthfulQA, CITA improves by \textbf{+0.054} (adjusting uncertainty appropriately), while DPO shows minimal change (\textbf{+0.001}). This validates CITA's ability to internalize calibration instructions.

    \item \textbf{Mandatory KL is essential}: The KL term---optional in DPO---becomes mandatory in CITA to prevent mode collapse and enable stable instruction-conditioned behavioral switching.

    \item \textbf{Instruction-alignment $\neq$ instruction-following}: Through the Instruction Switch Dataset, we demonstrate that CITA achieves \textit{deeper} instruction-alignment (internalized policies) rather than \textit{superficial} instruction-following (surface compliance).
\end{enumerate}

\vspace{-2mm}
\subsection{Implications}
\label{sec:implications}

CITA enables a new paradigm of \textbf{interactive and adaptive AI systems} where:

\begin{itemize}[leftmargin=*,itemsep=2pt]
    \item \textbf{Policy updates via natural language}: Alignment changes can be expressed through instructions without model retraining.

    \item \textbf{Context-aware responses}: Different users, applications, or deployment contexts can receive appropriately calibrated behaviors.

    \item \textbf{Dynamic safety policies}: Safety constraints can be adjusted at inference time based on deployment requirements (e.g., stricter for child-facing apps, more permissive for research contexts).

    \item \textbf{Reduced alignment tax}: No need for multiple specialized models for different alignment policies---a single CITA model can switch behaviors on demand.
\end{itemize}

\textbf{Outlook.} CITA opens several research directions: (1) scaling to larger models~\cite{touvron2023llama,jiang2023mistral,team2023gemini,chowdhery2022palm} and more instruction types, (2) theoretical analysis of instruction-conditioned learning dynamics~\cite{azar2023general}, (3) applications to multi-stakeholder alignment where different parties have different policy requirements~\cite{sun2024trustllm}, and (4) integration with constitutional AI methods~\cite{bai2022constitutional} for hierarchical instruction handling.

\vspace{-2mm}
\subsection{Reproducibility}
\label{sec:reproducibility}

All code, trained models, and the Instruction Switch Dataset are publicly available:

\begin{itemize}[leftmargin=*,itemsep=1pt]
    \item \textbf{Code}: \url{https://github.com/kapilw25/finetuning_evaluation}
    \item \textbf{Dataset}: \url{https://huggingface.co/datasets/kapilw25/ISD-Instruction-Switch-Dataset}
\end{itemize}
